\documentclass[header={Lecture 1 Exam Study Notes}]{study-note}

\newcommand{\Lp}{L^p}

\begin{document}

\StudyTitleBlock{01}{Lecture 1: Exam Study Notes}{Introduction, norms, metrics, completeness, continuity, and first operator definitions}

\vspace{0.8em}

\begin{focusbox}
This note is an \textbf{exam-focused rewrite} of the first handwritten lecture, filtered using the syllabus, theorem list, and oral-exam documents. It keeps the material you are expected to know and omits non-exam narrative such as historical remarks.
\end{focusbox}

\tableofcontents

\section{What To Learn From This Lecture}

\begin{focusbox}
For the exam, the core takeaways from Lecture 1 are:
\begin{itemize}
\item what functional analysis studies and why functionals matter;
\item the basic ambient spaces: vector, pre-Hilbert, normed, metric, and topological spaces;
\item convergence, Cauchy sequences, completeness, Banach and Hilbert spaces;
\item standard examples: $\R^d$, $\C^d$, $C^k(\Omega)$, $\ell^p$, $\Lp(\Omega)$, $W^{k,p}(\Omega)$;
\item Theorems \textbf{2.1--2.4} from the exam theorem list;
\item the first definitions for linear maps and bounded linear operators.
\end{itemize}
\end{focusbox}

\section{Functional Analysis In One Line}

Functional analysis studies mappings
\[
L:X\to Y,
\]
where $X$ and $Y$ are vector spaces over $\K$, with $\K=\R$ or $\C$, and where the spaces usually carry extra structure that allows one to discuss continuity, convergence, compactness, or solvability.

\begin{defbox}{Functional and ambient structures}
An important special case is $Y=\K$. Then
\[
L:X\to \K
\]
is called a \textbf{functional}.

Typical structures on $X$ are:
\begin{itemize}
\item a scalar product $(\cdot,\cdot)$: pre-Hilbert / inner-product space;
\item a norm $\norm{\cdot}$: normed space;
\item a distance $d$: metric space;
\item a topology $\mathcal T$: topological space.
\end{itemize}
\end{defbox}

\section{Core Examples}

\begin{exbox}{Examples you should recognize immediately}
\begin{itemize}
\item $\R^d$ and $\C^d$ with the usual Euclidean norm are finite-dimensional Banach spaces; with the standard inner product they are Hilbert spaces.
\item $C(\ol\Omega)$ with the sup norm
\[
\norm{f}_{\infty}=\max_{x\in\ol\Omega}\abs{f(x)}
\]
is Banach.
\item $\ell^p$ for $1\le p\le \infty$ are Banach spaces; $\ell^2$ is Hilbert.
\item $\Lp(\Omega)$ for $1\le p\le \infty$ are Banach spaces; $L^2(\Omega)$ is Hilbert.
\item Sobolev spaces $W^{k,p}(\Omega)$ are Banach; when $p=2$, one writes
\[
H^k(\Omega):=W^{k,2}(\Omega),
\]
and these are Hilbert spaces.
\end{itemize}
\end{exbox}

\begin{warningbox}{Finite vs. infinite dimensional intuition}
In $\R^d$, the span of the canonical basis is the whole space. In contrast, in $\ell^p$ the span of $e_1,e_2,\dots$ consists only of finitely supported sequences, so it is \emph{not} all of $\ell^p$.
\end{warningbox}

\section{Norms And Metrics}

\begin{defbox}{Norm}
Let $X$ be a vector space over $\K$. A map $x\mapsto \norm{x}_X$ is a \textbf{norm} if for all $x,y\in X$ and $\lambda\in\K$:
\begin{align*}
\text{(N1)}\quad & \norm{x}_X=0 \iff x=0,
\\
\text{(N2)}\quad & \norm{\lambda x}_X=\abs{\lambda}\norm{x}_X,
\\
\text{(N3)}\quad & \norm{x+y}_X\le \norm{x}_X+\norm{y}_X.
\end{align*}
Then $(X,\norm{\cdot}_X)$ is a \textbf{normed space}.
\end{defbox}

\begin{thmbox}{Theorem 2.1: A norm generates a metric}
\thmNormInducesMetric
\end{thmbox}

\begin{prooflevelbox}{FULL}
Know how to derive each metric property from the corresponding norm axiom. The key step is: triangle inequality for $d$ is exactly (N3).
\end{prooflevelbox}

\paragraph{Proof.}
Non-negativity and the implication $d(x,y)=0 \iff x=y$ follow from (N1). Symmetry follows from
\[
d(x,y)=\norm{x-y}_X=\norm{(-1)(y-x)}_X=\abs{-1}\norm{y-x}_X=d(y,x).
\]
The triangle inequality is exactly (N3):
\[
d(x,z)=\norm{x-z}_X=\norm{(x-y)+(y-z)}_X\le \norm{x-y}_X+\norm{y-z}_X=d(x,y)+d(y,z).
\]

\begin{focusbox}
The lecture also stresses that the metric induced by a norm is compatible with vector-space structure:
\[
d(x+z,y+z)=d(x,y),
\qquad
d(\lambda x,\lambda y)=\abs{\lambda}d(x,y),
\]
and norm balls are convex.
\end{focusbox}

\section{Convergence, Completeness, And Closed Subspaces}

\begin{defbox}{Convergence and Cauchy sequence}
In a normed space $(X,\norm{\cdot}_X)$:
\begin{itemize}
\item $x_n\to x$ means $\norm{x_n-x}_X\to 0$;
\item $(x_n)$ is \textbf{Cauchy} if
\[
\forall \varepsilon>0\ \exists n_0\ \forall m,n\ge n_0:\ \norm{x_n-x_m}_X<\varepsilon.
\]
\end{itemize}
\end{defbox}

\begin{defbox}{Completeness}
\begin{itemize}
\item A normed space is \textbf{Banach} if every Cauchy sequence converges in the space.
\item A scalar-product space is \textbf{Hilbert} if it is complete in the induced norm.
\item More generally, a metric linear space built from a separating sequence of seminorms and complete for that metric is called a \textbf{Fr\'echet space}. In the course this becomes Theorem~5.1 later.
\end{itemize}
\end{defbox}

\begin{thmbox}{Theorem 2.2: Characterization of closure and closed sets in a metric space}
\thmClosureSequential
\end{thmbox}

\begin{prooflevelbox}{FULL}
Both directions are short: $\bar A$ meets every ball around $x$, so pick $x_n\in A\cap B_{1/n}(x)$. Conversely, if $x_n\in A\to x$, every neighborhood of $x$ contains some $x_n$.
\end{prooflevelbox}

\paragraph{Proof.}
If $x\in\ol A$, then every ball $B_{1/n}(x)$ meets $A$; choose $x_n\in A\cap B_{1/n}(x)$. Then $x_n\to x$. Conversely, if $x_n\in A$ and $x_n\to x$, every neighborhood of $x$ contains some $x_n$, so $x\in\ol A$.

For the second claim, if $A$ is closed and $x_n\in A$ converges to $x$, then $x\in A=\ol A$. Conversely, if every convergent sequence in $A$ stays in $A$, then the first part shows that every point of $\ol A$ already lies in $A$, hence $A$ is closed.

\begin{thmbox}{Theorem 2.3: A subspace of a Banach space is complete iff it is closed}
\thmSubspaceCompleteIffClosed
\end{thmbox}

\begin{prooflevelbox}{FULL}
Know both directions: complete $\Rightarrow$ closed (Cauchy sequence converges in the subspace, limit must agree with the limit in $X$), and closed $\Rightarrow$ complete (Cauchy in $Y$ converges in $X$, closedness forces the limit back into $Y$).
\end{prooflevelbox}

\paragraph{Proof.}
If $Y$ is complete and $y_n\in Y$ converges in $X$ to some $x$, then $(y_n)$ is Cauchy in $Y$, hence converges in $Y$ to a point of $Y$; uniqueness of limits gives $x\in Y$. So $Y$ is closed.

Conversely, if $Y$ is closed and $(y_n)$ is Cauchy in $Y$, then it is Cauchy in $X$. Since $X$ is Banach, $y_n\to x$ for some $x\in X$. Closedness of $Y$ gives $x\in Y$, hence $Y$ is complete.

\section{Continuity And Topology}

\begin{defbox}{Continuity at a point}
Let $(X,\norm{\cdot}_X)$ and $(Y,\norm{\cdot}_Y)$ be normed spaces. A map $f:X\to Y$ is \textbf{continuous at} $x_0\in X$ if
\[
\forall \varepsilon>0\ \exists \delta>0\ \forall x\in X:\
\norm{x-x_0}_X<\delta \implies \norm{f(x)-f(x_0)}_Y<\varepsilon.
\]
In metric spaces this is equivalent to sequential continuity:
\[
x_n\to x_0 \implies f(x_n)\to f(x_0).
\]
\end{defbox}

\begin{thmbox}{Theorem 2.4: Continuity via preimages of open sets}
\thmContinuityPreimages
\end{thmbox}

\begin{prooflevelbox}{FULL}
Both directions are standard: continuity $\Rightarrow$ preimages of open sets are open (use the $\varepsilon$-$\delta$ definition), and the converse (take $U$ a neighborhood of $f(x)$, then $f^{-1}(U)$ is an open neighborhood of $x$).
\end{prooflevelbox}

\paragraph{Proof.}
If $f$ is continuous and $U$ is open in $Y$, then every point of $f^{-1}(U)$ has a neighborhood mapped into $U$, so $f^{-1}(U)$ is open. Conversely, if preimages of open sets are open, then for every neighborhood $U$ of $f(x)$, the set $f^{-1}(U)$ is an open neighborhood of $x$, which is exactly continuity at $x$.

\section{First Definitions About Linear Operators}

\begin{defbox}{Linear mapping}
A map
\[
L:\operatorname{Dom}(L)\subset X\to Y
\]
is \textbf{linear} if for all $x,y\in \operatorname{Dom}(L)$ and all $\alpha\in\K$:
\[
L(x+y)=Lx+Ly,
\qquad
L(\alpha x)=\alpha Lx.
\]
The domain $\operatorname{Dom}(L)$ is a subspace of $X$. We also use:
\[
\operatorname{Range}(L)=\operatorname{Im}(L),
\qquad
\ker L=\operatorname{Null}(L).
\]
\end{defbox}

\begin{defbox}{Bounded linear operator and operator norm}
If $L:X\to Y$ is linear on all of $X$, we call it \textbf{bounded} if
\[
\norm{L}:=\sup_{\norm{x}_X\le 1}\norm{Lx}_Y<\infty.
\]
\end{defbox}

\begin{thmbox}{Basic estimate for bounded linear operators}
If $L:X\to Y$ is linear and bounded, then
\[
\norm{Lx}_Y\le \norm{L}\,\norm{x}_X
\qquad\text{for all }x\in X.
\]
Hence bounded subsets of $X$ are mapped into bounded subsets of $Y$.
\end{thmbox}

\paragraph{Proof.}
If $x=0$, the claim is trivial. For $x\neq 0$,
\[
\norm{Lx}_Y
=\norm{x}_X\left\|L\left(\frac{x}{\norm{x}_X}\right)\right\|_Y
\le \norm{x}_X \sup_{\norm{u}_X\le 1}\norm{Lu}_Y
=\norm{L}\,\norm{x}_X.
\]

\begin{warningbox}{Exam orientation}
The \emph{major} operator theorems
\begin{itemize}
\item boundedness $\Leftrightarrow$ continuity for linear maps,
\item $\mathcal L(X,Y)$ is Banach when $Y$ is Banach,
\item finite-dimensional consequences,
\end{itemize}
belong mainly to the next lecture block, even though this first handwritten note already starts defining them.
\end{warningbox}

\section{What To Memorize Precisely}

\begin{focusbox}
Before the exam, you should be able to state without hesitation:
\begin{enumerate}
\item the definitions of normed space, metric space, Banach space, Hilbert space, continuity, linear map, bounded linear operator;
\item Theorems \textbf{2.1--2.4};
\item the standard examples $\R^d$, $C(\ol\Omega)$, $\ell^p$, $\Lp(\Omega)$, $W^{k,p}(\Omega)$;
\item the operator norm formula
\[
\norm{L}=\sup_{\norm{x}\le 1}\norm{Lx}
\]
and the estimate $\norm{Lx}\le \norm{L}\norm{x}$;
\item the finite-dimensional vs. infinite-dimensional contrast in sequence spaces.
\end{enumerate}
\end{focusbox}

\end{document}
